\chapter{Chapter05 — FKN Improvement}
%%% generated by the blueprint pipeline from FABL.Chapter05.FKNImprovement
\section{Overview}

Declarations of the Lean module \texttt{FABL.Chapter05.FKNImprovement}.

\section{Declarations}

\begin{definition}[Balanced FKN lift]
\lean{FABL.balancedFKNLift}
\statementsource{Analysis-of-Boolean-Functions-by-Ryan-ODonnell}{
  pdf-ffbe9c7b3306-p176-b002,
  pdf-ffbe9c7b3306-p252-b002
}
\label{FABL.balancedFKNLift}
\leanok
\uses{FABL.BooleanFunction}
For a Boolean function $f$ on the $n$-dimensional sign cube, define a Boolean function
on the $(n+1)$-dimensional sign cube by sending $z$ to the product of $z_0$ and $f$
evaluated at the vector whose $i$th coordinate is $z_0z_{i+1}$, for $0\le i<n$.
\end{definition}

\begin{theorem}[Lift at a positive leading sign]
\lean{FABL.balancedFKNLift_fin_cons_one}
\label{FABL.balancedFKNLift_fin_cons_one}
\leanok
\uses{FABL.BooleanFunction, FABL.balancedFKNLift}
\difficulty{1}
Let $n$ be a natural number, let $f$ be a function from $\\{-1,1\\}^n$ to the signs
$\\{\\pm1\\}$, and let $x\\in\\{-1,1\\}^n$. The balanced lift of $f$, evaluated at the
vector with first coordinate $1$ and remaining coordinates $x$, equals $f(x)$.
\end{theorem}

\begin{theorem}[Lift at a negative leading sign]
\lean{FABL.balancedFKNLift_fin_cons_neg_one}
\label{FABL.balancedFKNLift_fin_cons_neg_one}
\leanok
\uses{FABL.BooleanFunction, FABL.balancedFKNLift}
\difficulty{3}
Let $n$ be a natural number, let $f$ be a function from $\\{-1,1\\}^n$ to the signs
$\\{\\pm1\\}$, and let $x\\in\\{-1,1\\}^n$. The balanced lift of $f$, evaluated at the
vector with first coordinate $-1$ and remaining coordinates $x$, equals $-f(-x)$.
\end{theorem}

\begin{theorem}[Real lift at a positive leading sign]
\lean{FABL.balancedFKNLift_toReal_fin_cons_one}
\label{FABL.balancedFKNLift_toReal_fin_cons_one}
\leanok
\uses{FABL.BooleanFunction, FABL.BooleanFunction.toReal, FABL.balancedFKNLift, FABL.balancedFKNLift_fin_cons_one}
\difficulty{1}
Let $n$ be a natural number, let $f$ be a function from $\\{-1,1\\}^n$ to the signs
$\\{\\pm1\\}$, and let $x\\in\\{-1,1\\}^n$. After interpreting the outputs as real
numbers, the balanced lift of $f$ at the vector with first coordinate $1$ and remaining
coordinates $x$ equals $f(x)$.
\end{theorem}

\begin{theorem}[Real lift at a negative leading sign]
\lean{FABL.balancedFKNLift_toReal_fin_cons_neg_one}
\label{FABL.balancedFKNLift_toReal_fin_cons_neg_one}
\leanok
\uses{FABL.BooleanFunction, FABL.BooleanFunction.toReal, FABL.balancedFKNLift, FABL.balancedFKNLift_fin_cons_neg_one, FABL.signValue}
\difficulty{4}
Let $n$ be a natural number, let $f$ be a function from $\\{-1,1\\}^n$ to the signs
$\\{\\pm1\\}$, and let $x\\in\\{-1,1\\}^n$. After interpreting the outputs as real
numbers, the balanced lift of $f$ at the vector with first coordinate $-1$ and remaining
coordinates $x$ equals $-f(-x)$.
\end{theorem}

\begin{theorem}[Mean of the balanced lift]
\lean{FABL.mean_balancedFKNLift}
\label{FABL.mean_balancedFKNLift}
\leanok
\uses{FABL.BooleanFunction, FABL.BooleanFunction.toReal, FABL.balancedFKNLift, FABL.balancedFKNLift_toReal_fin_cons_neg_one, FABL.balancedFKNLift_toReal_fin_cons_one, FABL.expect_fin_cons, FABL.mean}
\difficulty{4}
For every Boolean function $f:\{-1,1\}^n\to\{-1,1\}$, define its lift on the
$(n+1)$-dimensional sign cube by $F(z)=z_1 f(z_1z_2,\ldots,z_1z_{n+1})$. The uniform
mean of the real-valued sign representation of $F$ is zero.
\end{theorem}

\begin{theorem}[Fourier coefficients under input negation]
\lean{FABL.fourierCoeff_negatedInput_toReal}
\label{FABL.fourierCoeff_negatedInput_toReal}
\leanok
\uses{FABL.BooleanFunction, FABL.BooleanFunction.toReal, FABL.fourierCoeff, FABL.monomial, FABL.signValue}
\difficulty{5}
Let $n$ be a natural number, let $f$ map the $n$-dimensional sign cube to $\{\pm1\}$,
and let $S$ be a set of coordinates. The Fourier coefficient at $S$ of the function
$x\mapsto -f(-x)$ is $-(-1)^{|S|}\widehat f(S)$.
\end{theorem}

\begin{theorem}[First Fourier coefficient of the balanced lift]
\lean{FABL.fourierCoeff_balancedFKNLift_zero}
\label{FABL.fourierCoeff_balancedFKNLift_zero}
\leanok
\uses{FABL.BooleanFunction, FABL.BooleanFunction.toReal, FABL.balancedFKNLift, FABL.balancedFKNLift_toReal_fin_cons_neg_one, FABL.balancedFKNLift_toReal_fin_cons_one, FABL.firstCoordinateSlice, FABL.fourierCoeff, FABL.fourierCoeff_insert_zero_tailFrequency, FABL.fourierCoeff_negatedInput_toReal, FABL.tailFrequency}
\difficulty{4}
For every natural number $n$ and every function $f:\{-1,1\}^n\to\{\pm1\}$, the Fourier
coefficient of the balanced lift of $f$ at the singleton frequency consisting of the new
coordinate equals the constant Fourier coefficient of $f$.
\end{theorem}

\begin{theorem}[Singleton Fourier coefficient under the balanced lift]
\lean{FABL.fourierCoeff_balancedFKNLift_succ}
\label{FABL.fourierCoeff_balancedFKNLift_succ}
\leanok
\uses{FABL.BooleanFunction, FABL.BooleanFunction.toReal, FABL.balancedFKNLift, FABL.balancedFKNLift_toReal_fin_cons_neg_one, FABL.balancedFKNLift_toReal_fin_cons_one, FABL.firstCoordinateSlice, FABL.fourierCoeff, FABL.fourierCoeff_negatedInput_toReal, FABL.fourierCoeff_tailFrequency, FABL.tailFrequency}
\difficulty{4}
Let $n$ be a nonnegative integer, let $f:\{-1,1\}^n\to\{-1,1\}$, and let
$i\in\{0,\ldots,n-1\}$. Define a function on $\{-1,1\}^{n+1}$ by $F(z)=z_0
f((z_0z_1,\ldots,z_0z_n))$. Then the Fourier coefficient of $F$ at the singleton
coordinate $i+1$ equals the Fourier coefficient of $f$ at coordinate $i$.
\end{theorem}

\begin{theorem}[Level-one weight of the balanced lift]
\lean{FABL.fourierWeightAtLevel_one_balancedFKNLift}
\label{FABL.fourierWeightAtLevel_one_balancedFKNLift}
\leanok
\uses{FABL.BooleanFunction, FABL.BooleanFunction.toReal, FABL.balancedFKNLift, FABL.fourierCoeff_balancedFKNLift_succ, FABL.fourierCoeff_balancedFKNLift_zero, FABL.fourierWeightAtLevel, FABL.fourierWeightAtLevel_one_eq_sum_singleton, FABL.fourierWeightAtMost, FABL.fourierWeightAtMost_one_eq_empty_add_sum_singleton}
\difficulty{3}
Let $f:\{-1,1\}^n\to\{-1,1\}$, and define its balanced lift on $\{-1,1\}^{n+1}$ by
$g(z)=z_0 f((z_0z_{i+1})_{0\le i<n})$. The level-one Fourier weight of $g$ equals the
sum of the squared Fourier coefficients of $f$ over all subsets of size at most one,
with coefficients computed under the uniform measure.
\end{theorem}

\begin{definition}[Balanced FKN projection]
\lean{FABL.balancedFKNProjection}
\statementsource{Analysis-of-Boolean-Functions-by-Ryan-ODonnell}{pdf-ffbe9c7b3306-p127-b002}
\label{FABL.balancedFKNProjection}
\leanok
\uses{FABL.BooleanFunction, FABL.signedDictator}
For a nonnegative integer $n$, an index $i\in\{0,\ldots,n\}$, and a choice of whether to
negate, define a Boolean function on $\{\pm1\}^n$ by evaluating the signed dictator at
$i$ on the vector obtained by prepending $+1$ to the input.
\end{definition}

\begin{theorem}[One-coordinate dependence of the projection]
\lean{FABL.balancedFKNProjection_isKJunta_one}
\label{FABL.balancedFKNProjection_isKJunta_one}
\leanok
\uses{FABL.IsKJunta, FABL.balancedFKNProjection, FABL.dictator, FABL.signedDictator}
\difficulty{4}
Let $n$ be a natural number, choose a coordinate $i$ of the $(n+1)$-dimensional sign
cube, and choose whether to negate its sign. The function on the $n$-dimensional sign
cube obtained by evaluating that signed coordinate function on vectors whose first
coordinate is $1$ depends on at most one coordinate.
\end{theorem}

\begin{theorem}[Signed dictator under global negation]
\lean{FABL.signedDictator_neg_input}
\label{FABL.signedDictator_neg_input}
\leanok
\uses{FABL.dictator, FABL.signedDictator}
\difficulty{1}
Let $n\ge 0$, let $i$ be a coordinate in $\{0,\ldots,n\}$, let
$\varepsilon\in\{\text{false},\text{true}\}$ specify whether the dictator is negated,
and let $z\in\{-1,1\}^{n+1}$. Negating every coordinate of $z$ negates the value of the
signed dictator at coordinate $i$.
\end{theorem}

\begin{theorem}[Signed dictator under a coordinate extension]
\lean{FABL.signedDictator_fin_cons_neg_one}
\label{FABL.signedDictator_fin_cons_neg_one}
\leanok
\uses{FABL.signedDictator, FABL.signedDictator_neg_input}
\difficulty{4}
Let $n\ge 0$, let $i$ be a coordinate in $\{0,\ldots,n\}$, let
$\varepsilon\in\{\text{false},\text{true}\}$ specify whether the dictator is negated,
and let $x\in\{-1,1\}^n$. Extend $x$ to length $n+1$ by placing $-1$ first. The signed
dictator at this extended vector equals the negative of its value at the vector formed
by placing $1$ first and negating every coordinate of $x$.
\end{theorem}

\begin{theorem}[Distance under the balanced lift]
\lean{FABL.relativeHammingDist_balancedFKNLift_signedDictator}
\label{FABL.relativeHammingDist_balancedFKNLift_signedDictator}
\leanok
\uses{FABL.BooleanFunction, FABL.balancedFKNLift, FABL.balancedFKNLift_fin_cons_neg_one, FABL.balancedFKNProjection, FABL.expect_fin_cons, FABL.relativeHammingDist, FABL.signedDictator, FABL.signedDictator_fin_cons_neg_one, FABL.uniformProbability, FABL.uniformProbability_ne_eq_relativeHammingDist}
\difficulty{5}
Let $f$ be a Boolean function on the sign cube of dimension $n$, let $i$ be a coordinate
of the sign cube of dimension $n+1$, and choose whether to negate the dictator on
coordinate $i$. The relative Hamming distance from the balanced lift of $f$ to that
signed dictator equals the relative Hamming distance from $f$ to the function obtained
by evaluating the signed dictator at the input formed by adjoining the sign $+1$ as its
first coordinate.
\end{theorem}

\begin{theorem}[Exists is k junta one relative hamming dist le of fourier weight at most one]
\lean{FABL.exists_isKJunta_one_relativeHammingDist_le_of_fourierWeightAtMost_one}
\statementsource{Analysis-of-Boolean-Functions-by-Ryan-ODonnell}{
  pdf-ffbe9c7b3306-p061-b001,
  pdf-ffbe9c7b3306-p252-b002
}
\label{FABL.exists_isKJunta_one_relativeHammingDist_le_of_fourierWeightAtMost_one}
\leanok
\uses{FABL.BooleanFunction, FABL.BooleanFunction.toReal, FABL.IsKJunta, FABL.balancedFKNLift, FABL.balancedFKNProjection, FABL.balancedFKNProjection_isKJunta_one, FABL.fkn, FABL.fourierWeightAtLevel, FABL.fourierWeightAtLevel_one_balancedFKNLift, FABL.fourierWeightAtMost, FABL.relativeHammingDist, FABL.relativeHammingDist_balancedFKNLift_signedDictator}
\difficulty{4}
% dataset-pipeline: stub -- prose pending build_informal_statements.py
Statement pending regeneration.
\end{theorem}

\begin{definition}[FKN improvement parameter]
\lean{FABL.fknImprovementEta}
\statementsource{Analysis-of-Boolean-Functions-by-Ryan-ODonnell}{
  pdf-ffbe9c7b3306-p061-b002,
  pdf-ffbe9c7b3306-p129-b005
}
\label{FABL.fknImprovementEta}
\leanok
For real numbers $C$ and $\delta$, define the parameter
$\eta=16C^2\delta^2\max\{\log_2(1/(C\delta)),1\}$.
\end{definition}

\begin{theorem}[A logarithmic inequality]
\lean{FABL.twentyFiveDivEight_lt_logb_two_nine}
\label{FABL.twentyFiveDivEight_lt_logb_two_nine}
\leanok
\difficulty{4}
The rational number $25/8$ is strictly less than the logarithm of $9$ to base $2$.
\end{theorem}

\begin{theorem}[Base-two logarithm of eight]
\lean{FABL.logb_two_eight}
\label{FABL.logb_two_eight}
\leanok
\difficulty{3}
The base-two logarithm of $8$ equals $3$.
\end{theorem}

\begin{theorem}[Base-two logarithm of four]
\lean{FABL.logb_two_four}
\label{FABL.logb_two_four}
\leanok
\difficulty{3}
The base-two logarithm of $4$ equals $2$.
\end{theorem}

\begin{theorem}[Logarithmic bound on a small parameter]
\lean{FABL.mul_logb_two_inv_lt_one_third}
\label{FABL.mul_logb_two_inv_lt_one_third}
\leanok
\difficulty{6}
For every real number $x$ with $0<x\le1/10$, one has $x\log_2(1/x)<1/3$.
\end{theorem}

\begin{theorem}[Large-parameter negativity]
\lean{FABL.exercise5_38a}
\label{FABL.exercise5_38a}
\leanok
\uses{FABL.fknImprovementEta, FABL.logb_two_eight, FABL.logb_two_four, FABL.twentyFiveDivEight_lt_logb_two_nine}
\difficulty{6}
Let $C,\delta\in\mathbb R$ satisfy $1\le C$, $0<\delta$, and $1/(10C)<\delta$. Then
$1-\delta/2-2\eta(C,\delta)<0$, where
$\eta(C,\delta)=16C^2\delta^2\max\{\log_2(1/(C\delta)),1\}$.
\end{theorem}

\begin{theorem}[Squared lower-bound comparison]
\lean{FABL.exercise5_38b}
\label{FABL.exercise5_38b}
\leanok
\uses{FABL.fknImprovementEta, FABL.mul_logb_two_inv_lt_one_third}
\difficulty{5}
Let $C,\delta\in\mathbb R$ satisfy $1\le C$, $0<\delta$, and $\delta\le 1/(10C)$. Then
$1-\delta-16C^2\delta^2\log_2(1/(C\delta))\ge (1-\delta/2-2\eta(C,\delta))^2$, where
$\eta(C,\delta)=16C^2\delta^2\max\{\log_2(1/(C\delta)),1\}$.
\end{theorem}

\begin{theorem}[Nonnegative lower bound from a squared estimate]
\lean{FABL.exercise5_38_nonnegative_lower_bound}
\label{FABL.exercise5_38_nonnegative_lower_bound}
\leanok
\uses{FABL.exercise5_38b, FABL.fknImprovementEta}
\difficulty{4}
Let $C,\delta,z\in\mathbb R$ satisfy $C\ge1$, $\delta>0$, $\delta\le1/(10C)$, and
$z\ge0$. If
\[
1-\delta-16C^2\delta^2\log_2\!\left(\frac1{C\delta}\right)\le z^2,
\]
then, defining
$\eta=16C^2\delta^2\max\!\left\{\log_2\!\left(\frac1{C\delta}\right),1\right\}$, one has
\[
1-\frac{\delta}{2}-2\eta\le z.
\]
\end{theorem}

\begin{theorem}[Square-root lower bound]
\lean{FABL.exercise5_38_sqrt_lower_bound}
\label{FABL.exercise5_38_sqrt_lower_bound}
\leanok
\uses{FABL.exercise5_38_nonnegative_lower_bound, FABL.fknImprovementEta}
\difficulty{3}
Let $C,\delta,A\in\mathbb R$ satisfy $1\le C$, $0<\delta$, $\delta\le 1/(10C)$, and
$0\le A$. If $1-\delta-16C^2\delta^2\log_2(1/(C\delta))\le A$, then
$1-\delta/2-2\eta(C,\delta)\le\sqrt A$, where
$\eta(C,\delta)=16C^2\delta^2\max\{\log_2(1/(C\delta)),1\}$.
\end{theorem}
